\documentclass[10pt,a4paper]{amsart}
\usepackage[margin=19mm]{geometry}
\usepackage{amsmath,amssymb,mathtools}
\usepackage[scr=rsfs,cal=euler]{mathalfa}
\usepackage{xfrac}
\usepackage[unicode,hidelinks,bookmarks=false]{hyperref}
\newcommand{\ie}{i.e.\ }
\newcommand{\cf}{cf.\ }
\newcommand{\resp}{respectively\ }
\newcommand{\Subsection}{\subsection}
\providecommand{\nobreakdash}{\nobreak\hbox{-}}
\setlength{\emergencystretch}{2em}
\multlinegap0pt
\usepackage{bm}
\usepackage{bbm}
\usepackage{dsfont}
\usepackage{xspace}
\usepackage{cancel}
\usepackage{mathtools}
\usepackage{amsthm}
\makeatletter
\@ifundefined{theorem}{\newtheorem{theorem}{Theorem}}{}
\makeatother
\title[Starfield: Demand-Aware Satellite Topology Design for Low-Ea]{Starfield: Demand-Aware Satellite Topology Design for Low-Earth Orbit Mega Constellations}
\author{Shayan Hamidi Dehshali, Tzu-Hsuan Liao, Shaileshh Bojja Venkatakrishnan}
\date{Source: arXiv:2601.10083v2 (CC BY 4.0)}
\begin{document}
\maketitle

\noindent\emph{Source and licence.} Starfield: Demand-Aware Satellite Topology Design for Low-Earth Orbit Mega Constellations. Version arXiv:2601.10083v2 (CC BY 4.0), distributed under the Creative Commons Attribution 4.0 International licence, as recorded in the arXiv licence metadata for this exact version and shown on the abstract page https://arxiv.org/abs/2601.10083v2. Full rights notice: licenses/arxiv-2601.10083-CC-BY-4.0.txt.\par\smallskip

\noindent\textit{Editorial scope.} Counted unit: the Theorem whose source label is \texttt{None} in the article (source0.tex lines 665--672, source label \texttt{None}), delivered together with the article's own complete original proof of it (source0.tex lines 673--694, one \texttt{proof} environment of 22 lines). No mathematics is written, repaired or re-derived: statement and proof are the article's own text line for line. Required context retained uncounted: none. Every definition, display and label of those lines is retained unchanged. Omitted from this package and not redistributed: 1--664, 695--812. Nothing inside a retained excerpt is omitted or altered: every retained body is delivered exactly as the article's lines are written, so no omission receipt of any kind accompanies this package. Citations: the retained excerpts carry no citation command at all, so no bibliography entry of the article is transcribed and no reference is redistributed. Reference adaptation: none was needed and none was made; every reference inside the delivered unit resolves inside it, so no unresolved reference or citation is produced. Printed slips kept literal: no printed word, symbol, hypothesis or number of the counted statement or of its proof was altered. Layout and class: the article's own class is \texttt{acmart} with packages including siunitx, algorithmicx, algpseudocode, caption, booktabs, mdframed, hyperref; the wrapper is the lane's pinned \texttt{amsart} editorial wrapper at 10pt on A4, which restates the theorem-like environments the retained text needs and adds the article's own macro definitions that the retained text needs (none), with extra wrapper packages (bm, bbm, dsfont, xspace, cancel, mathtools). The delivered numbering is the unit's own. Assets: no retained span contains a figure environment, a graphics-inclusion command, a PDF-page inclusion, a table or a TikZ picture; no image file is copied into this package, the package declares no asset dependency and no image file is redistributed with it. Licence: version arXiv:2601.10083v2 of this article is distributed under the Creative Commons Attribution 4.0 International licence, as recorded in the arXiv licence metadata for this exact version and shown on the abstract page https://arxiv.org/abs/2601.10083v2. A full rights notice is delivered at licenses/arxiv-2601.10083-CC-BY-4.0.txt. Characters: every retained excerpt is pure ASCII, so the delivered engine, XeLaTeX with its default Latin Modern text font, reports no missing character.
\begin{theorem}
Consider any topology $E$. 
For any continuous, compact region $\mathcal{A} \subset \mathbb{R} \times \mathbb{R}$ with a non-empty $\mathcal{D}_\mathcal{A}$ and parameter $\epsilon \in [0, \pi)$, there exists a demand $(s,d) \in \mathcal{D}_\mathcal{A}$ with a stretch 
\begin{align}
\sigma_E(s,d) \geq \frac{ \min \left(l(s,d), \frac{R \lambda_\mathcal{A} (2\alpha(\mathcal{A}) \rho^2 + 1) \delta}{ |\mathcal{D}_\mathcal{A}|} \right)  + \frac{l(s,d) -   \min \left(l(s,d), \frac{R \lambda_\mathcal{A} (2\alpha(\mathcal{A}) \rho^2 + 1) \delta}{ |\mathcal{D}_\mathcal{A}|} \right)}{| \cos(\epsilon) |} }{ l(s,d) },
\end{align} 
where $\lambda_\mathcal{A}$ is the maximum number of demands in $\mathcal{D}_\mathcal{A}$ that are oriented at an angle within $\epsilon$ of any reference direction.   
\end{theorem}

\begin{proof}
The minimum number of edges needed in $E_\mathcal{A}^{(s,d)}(\epsilon)$  to satisfy a demand $(s, d)$ is $\frac{l(s,d)}{ R }$.
The minimum occurs if all the edges in $E_\mathcal{A}^{(s,d)}(\epsilon)$ are aligned precisely with $(s, d)$.
If $|E_\mathcal{A}^{(s,d)}(\epsilon)| < \frac{l(s,d)}{R}$, it implies there is at least one edge in the $s$ to $d$ shortest path who orientation angle differs from $\theta(s,d)$ by more than $\epsilon$. 
If $|E_\mathcal{A}^{(s,d)}(\epsilon)| = 0$, it implies the stretch of $(s, d)$ is at least $\frac{1}{ |\cos(\epsilon)| }$.  
Note that 
\begin{align}
\sum_{(u, v) \in E_\mathcal{A}} | \mathcal{D}_\mathcal{A}^{(u,v)}(\epsilon) | = \sum_{(s, d) \in \hat{\mathcal{D}}_\mathcal{A}} | E_\mathcal{A}^{(s, d)}(\epsilon)|.
\end{align}   

Suppose the demands are such that for any orientation of an edge $(u, v)$ we have $|\mathcal{D}_\mathcal{A}^{(u,v)}(\epsilon)|$ is less than $\lambda_\mathcal{A}$. 
Then 
\begin{align}
\sum_{(u, v) \in E_\mathcal{A}} | \mathcal{D}_\mathcal{A}^{(u,v)}(\epsilon) | \leq \lambda_\mathcal{A} |E_\mathcal{A}| = \lambda_\mathcal{A} (2\alpha(\mathcal{A}) \rho^2 + 2) \delta.
\end{align} 
This implies  $\sum_{(s, d) \in \mathcal{D}_\mathcal{A}} | E_\mathcal{A}^{(s, d)}(\epsilon)| \leq \lambda_\mathcal{A} (2\alpha(\mathcal{A}) \rho^2 + 2) \delta$.  
Therefore, there exists at least one demand $(s, d) \in \mathcal{D}_\mathcal{A}$ such that $|E_\mathcal{A}^{(s,d)}(\epsilon)| \leq \frac{\lambda_\mathcal{A} (2\alpha(\mathcal{A}) \rho^2 + 2) \delta}{ |\mathcal{D}_\mathcal{A}| }$.
To lower bound the stretch for this demand, we consider the best case scenario in which all $\frac{\lambda_\mathcal{A} (2\alpha(\mathcal{A}) \rho^2 + 2) \delta}{ |\mathcal{D}_\mathcal{A}| }$ edges are aligned in the exact same direction as $(s,d)$ and the remaining edges on the path from $s$ to $d$ are aligned at an angle exactly $\epsilon$ away from $(s,d)$.  
The total length of the edges that are aligned along $(s,d)$ is $ \min \left(l(s,d), \frac{R \lambda_\mathcal{A} (2\alpha(\mathcal{A}) \rho^2 + 2) \delta}{ |\mathcal{D}_\mathcal{A}|} \right) $.  
The total length of the edges that are aligned at an angle $\epsilon$ away from $(s,d)$ is $\frac{l(s,d) -   \min \left(l(s,d), \frac{R \lambda_\mathcal{A} (2 \alpha(\mathcal{A}) \rho^2 + 2) \delta}{ |\mathcal{D}_\mathcal{A}|} \right)}{|\cos(\epsilon)|}$.  
This gives a lower bound on the stretch as mentioned in the theorem.  
\end{proof} 

\end{document}
